\documentclass[11pt,a4paper]{article}

\usepackage[utf8]{inputenc}
\usepackage[T1]{fontenc}
\usepackage{lmodern}
\usepackage[french]{babel}
\usepackage[margin=2.5cm]{geometry}
\usepackage{amsmath,amssymb}
\usepackage{booktabs}
\usepackage{array}
\usepackage{xcolor}
\usepackage{listings}
\usepackage{tikz}
\usepackage[french,ruled,vlined,linesnumbered]{algorithm2e}
\SetKwInput{KwIn}{Entrées}
\SetKwInput{KwOut}{Sortie}
\SetKwFor{For}{pour}{faire}{fin}
\SetKwIF{If}{ElseIf}{Else}{si}{alors}{sinon si}{sinon}{fin}
\SetKw{KwTo}{à}
\SetKw{KwRet}{retourner}
\usepackage[colorlinks=true,linkcolor=blue!50!black,urlcolor=blue!50!black,citecolor=blue!50!black]{hyperref}

\lstset{
  language=C++,
  basicstyle=\ttfamily\small,
  keywordstyle=\color{blue!60!black}\bfseries,
  commentstyle=\color{green!40!black}\itshape,
  stringstyle=\color{red!50!black},
  numbers=left,
  numberstyle=\tiny\color{gray},
  frame=single,
  breaklines=true,
  columns=fullflexible,
  keepspaces=true,
  showstringspaces=false
}

\newcommand{\code}[1]{\texttt{#1}}
\newcommand{\N}{\mathbb{N}}

\title{Optimisation de la hauteur de l'octree de la méthode multipolaire rapide\\
dans feeLLGood\\[0.5em]
\large Courbe de Morton, comptage des feuilles non vides et choix de la hauteur}
\author{Branche \code{optim\_tree\_height}}
\date{Novembre 2026}

\begin{document}
\sloppy
\maketitle

\begin{abstract}
Le champ démagnétisant de feeLLGood est calculé par la méthode multipolaire rapide (FMM) de
ScalFMM~1.5. Jusqu'ici, l'octree de la FMM avait une hauteur fixe de 6 niveaux, quel que soit le
maillage. Ce document explique pourquoi cette hauteur est presque toujours mauvaise, comment
choisir automatiquement la bonne hauteur à partir de la répartition des particules, et comment ce
choix a été mis en œuvre, à l'image de st-feeLLGood, à l'aide des clés de Morton. On détaille la
courbe de Morton, son lien avec l'octree, l'algorithme de comptage des feuilles non vides, le code
C++ correspondant, les tests unitaires, puis les mesures de performance sur un portable (8 threads)
et sur une machine de 96 cœurs (jusqu'à 192 threads). Le temps du calcul magnétostatique est divisé
par 2,6 à 18,5 selon le maillage et le nombre de threads, sans perte de précision.
\end{abstract}

\tableofcontents

%======================================================================
\section{Contexte}
%======================================================================

\subsection{Le calcul du champ démagnétisant}

Le champ démagnétisant dérive d'un potentiel scalaire $\phi$, $\mathbf{H}_d = -\nabla\phi$, créé
par les charges magnétiques volumiques $\rho = -M_s\,\nabla\cdot\mathbf{u}$ et surfaciques
$\sigma = M_s\,\mathbf{u}\cdot\mathbf{n}$ :
\begin{equation}
\phi(\mathbf{x}) = \frac{1}{4\pi}\left[\int_\Omega \frac{\rho(\mathbf{y})}{|\mathbf{x}-\mathbf{y}|}\,
\mathrm{d}V_y + \int_{\partial\Omega} \frac{\sigma(\mathbf{y})}{|\mathbf{x}-\mathbf{y}|}\,
\mathrm{d}S_y\right].
\label{eq:phi}
\end{equation}
Les intégrales sont discrétisées aux points de Gauss des éléments : chaque point de Gauss
$\mathbf{y}_g$ porte une charge ponctuelle $q_g$ (densité multipliée par le poids d'intégration).
Le potentiel aux nœuds du maillage est alors une somme de type $N$ corps,
\begin{equation}
\phi_i = \frac{1}{4\pi}\sum_g \frac{q_g}{|\mathbf{x}_i - \mathbf{y}_g|},
\label{eq:nbody}
\end{equation}
complétée par des corrections analytiques locales pour les charges surfaciques
(\code{Tri::correctionCharges}). Dans le code (\code{src/fmm\_demag.h}), on a :
\begin{itemize}
  \item des \emph{cibles} : les nœuds magnétiques, où l'on veut $\phi$ ;
  \item des \emph{sources} : les points de Gauss des tétraèdres magnétiques (\code{Tetra::NPI}
        $= 5$ points par tétraèdre) et des triangles magnétiques (\code{Triangle::NPI} $= 3$
        points par triangle).
\end{itemize}
Le nombre total de particules (cibles et sources) est donc
\begin{equation}
N = N_{\text{nœuds}} + 5\,N_{\text{tétra}} + 3\,N_{\text{tri}},
\end{equation}
soit environ 25 à 30 fois le nombre de nœuds. Par exemple, le maillage \code{example0.1}
(3990 nœuds, 19339 tétraèdres, 3370 triangles) compte $N = 110\,795$ particules.

La somme~\eqref{eq:nbody} coûte $O(N^2)$ si on l'évalue directement ; la FMM la calcule en
$O(N)$ avec une précision contrôlée. Elle est évaluée deux fois par pas de temps : une fois pour
le potentiel de $\mathbf{u}$ et une fois pour celui de $\mathbf{v}$ (option par défaut
\code{FIRST\_ORDER=OFF}, schéma d'ordre 2 en temps).

\subsection{Normalisation et boîte de calcul}

Avant insertion dans l'octree, les positions sont centrées sur le centre $\mathbf{c}$ de la boîte
englobante du maillage et mises à l'échelle :
\begin{equation}
\tilde{\mathbf{x}} = \frac{2}{\ell_{\max}}\,(\mathbf{x} - \mathbf{c}),
\qquad \ell_{\max} = \max(\ell_x, \ell_y, \ell_z),
\end{equation}
où $\ell_x, \ell_y, \ell_z$ sont les dimensions de la boîte englobante (\code{msh.l}). Toutes les
particules sont ainsi dans le cube $[-1,1]^3$. La boîte de l'octree est un cube de côté
$w = 2{,}01$ centré sur l'origine (constante \code{boxWidth}), légèrement plus grand pour éviter
les particules exactement sur le bord.

\subsection{La méthode multipolaire rapide et son octree}

La FMM range les particules dans un \emph{octree} : la boîte racine (niveau 0) est découpée en
8 cubes (niveau 1), chacun à nouveau en 8 (niveau 2), etc. Dans ScalFMM, un octree de
\emph{hauteur} $h$ possède les niveaux $0, 1, \dots, h-1$ ; les \emph{feuilles} sont les cellules
du niveau $h-1$, au nombre de $2^{h-1}$ par direction, soit $8^{h-1}$ au plus, de côté
\begin{equation}
w_{h} = \frac{w}{2^{h-1}}.
\end{equation}

Le potentiel en une cible est la somme de deux contributions :
\begin{itemize}
  \item le \textbf{champ proche} (P2P, \emph{particle to particle}) : interaction directe avec les
        particules de la feuille de la cible et de ses 26 feuilles voisines ;
  \item le \textbf{champ lointain}, calculé par développements en harmoniques sphériques :
        \begin{itemize}
          \item P2M : les charges d'une feuille sont résumées par un développement multipolaire ;
          \item M2M : remontée des développements multipolaires des enfants vers le parent ;
          \item M2L : conversion du développement multipolaire d'une cellule en développement
                local dans chaque cellule de sa \emph{liste d'interaction} : cellules du même
                niveau, enfants des voisins du parent, mais non voisines elles-mêmes. Il y en a
                au plus $6^3 - 3^3 = 189$ ;
          \item L2L : descente des développements locaux du parent vers les enfants ;
          \item L2P : évaluation du développement local aux cibles d'une feuille.
        \end{itemize}
\end{itemize}

feeLLGood utilise le noyau \og rotation\fg{} de ScalFMM (\code{FRotationKernel}), avec des
développements tronqués à l'ordre $P = 9$, soit $(P+1)(P+2)/2 = 55$ coefficients complexes par
cellule. Avec ce noyau, un opérateur M2L coûte $O(P^3)$ opérations (rotation, translation axiale,
rotation inverse).

Un point essentiel pour la suite : ScalFMM \emph{n'alloue que les cellules non vides}. À
l'insertion d'une particule, les cellules qui la contiennent sont créées si elles n'existent pas
encore (\code{FSubOctree::insert}) ; les itérations des différentes passes ne parcourent que les
cellules existantes. Le coût de la FMM dépend donc du nombre de cellules \emph{non vides}, et non
du nombre $8^{h-1}$ de cellules possibles.

%======================================================================
\section{Pourquoi la hauteur compte}
%======================================================================

\subsection{Le compromis champ proche / champ lointain}

Notons $L_h$ le nombre de feuilles non vides d'un octree de hauteur $h$ et
\begin{equation}
s_h = \frac{N}{L_h}
\end{equation}
le nombre moyen de particules par feuille non vide. Un modèle de coût simple, en supposant les
particules réparties de façon à peu près uniforme dans les feuilles non vides, est :
\begin{itemize}
  \item \textbf{champ proche} : chaque particule interagit avec les particules de 27 feuilles,
        d'où un coût $\approx \alpha \cdot 27\, N\, s_h$ ;
  \item \textbf{champ lointain} : dominé par les M2L, 189 au plus par cellule, d'où un coût
        $\approx \beta\, P^3 \cdot 189\, L_h$ pour le niveau des feuilles, les niveaux supérieurs
        ajoutant au plus $1/7$ de plus (série géométrique de raison $1/8$).
\end{itemize}
Avec $L_h = N/s_h$, le coût total s'écrit
\begin{equation}
T(s) \approx N \left( 27\,\alpha\, s + \frac{189\,\beta\, P^3}{s} \right),
\label{eq:cost}
\end{equation}
qui est minimal pour
\begin{equation}
s^\ast = \sqrt{\frac{7\,\beta\,P^3}{\alpha}}.
\label{eq:sopt}
\end{equation}

Ce modèle, même grossier, donne trois enseignements :
\begin{enumerate}
  \item \textbf{L'optimum ne dépend pas de $N$}, mais seulement du nombre moyen de particules par
        feuille. Une hauteur fixe ne peut donc pas convenir à tous les maillages : il faut une
        hauteur qui croît avec $N$, à peu près comme $\log_8 N$.
  \item \textbf{L'optimum dépend du noyau} ($P$, coût relatif $\beta/\alpha$ des M2L et des
        P2P), donc il faut le calibrer par la mesure.
  \item \textbf{Trop de niveaux coûte très cher} : quand $s$ devient petit, le terme en $1/s$
        explose. Les feuilles ne contiennent plus qu'une ou deux particules, et chaque
        particule paye jusqu'à 189 conversions M2L de 55 coefficients pour remplacer quelques
        interactions directes.
\end{enumerate}
En pratique, le modèle sous-estime l'asymétrie : les P2P sont des boucles simples, bien
vectorisées, alors que les M2L sont plus coûteux qu'estimé (accès mémoire, synchronisations entre
threads à chaque niveau). Les mesures montrent qu'un niveau de trop coûte bien plus qu'un niveau
de moins (section~\ref{sec:calib}).

\subsection{Le cas de la hauteur fixe à 6}

Avec $h = 6$, il y a jusqu'à $32^3 = 32\,768$ feuilles. Pour \code{example0.1}
($N = 110\,795$), même si toutes les feuilles étaient remplies, on aurait en moyenne 3,4
particules par feuille ; en réalité le maillage (un ellipsoïde) ne remplit pas toute la boîte, mais
l'ordre de grandeur reste de quelques particules par feuille. On est dans le régime où le champ
lointain domine largement. Les mesures de la section~\ref{sec:results} le confirment : la hauteur
adaptée ($h = 4$) est 4 fois plus rapide à 8 threads.

\subsection{Pourquoi compter les feuilles non vides}

Le nombre moyen de particules par feuille doit être calculé sur les feuilles non vides, puisque
ce sont les seules que ScalFMM crée et traite. Prenons deux maillages de même nombre de
particules : une sphère, qui remplit à peu près la moitié de la boîte, et un fil long, ou une
couche mince, qui n'en occupe qu'une petite fraction. À hauteur égale, le fil a beaucoup moins de
feuilles non vides, donc beaucoup plus de particules par feuille : il lui faut un octree plus
profond. Une formule du type $h = \lceil \log_8 (N/s^\ast) \rceil + 1$, qui supposerait la boîte
remplie, sous-estimerait la hauteur pour les géométries allongées ou plates. D'où la nécessité de
compter réellement les feuilles non vides, pour chaque hauteur candidate.

%======================================================================
\section{La courbe de Morton}
%======================================================================

\subsection{Définition}

Considérons une grille de $2^m \times 2^m \times 2^m$ cellules. Une cellule est repérée par ses
indices entiers $(i, j, k) \in \{0, \dots, 2^m-1\}^3$, écrits en binaire sur $m$ bits :
\begin{equation}
i = \sum_{b=0}^{m-1} i_b\, 2^b, \qquad
j = \sum_{b=0}^{m-1} j_b\, 2^b, \qquad
k = \sum_{b=0}^{m-1} k_b\, 2^b,
\qquad i_b, j_b, k_b \in \{0, 1\}.
\end{equation}
La \textbf{clé de Morton} de la cellule est l'entier de $3m$ bits obtenu en \emph{entrelaçant} les
bits des trois indices :
\begin{equation}
K(i, j, k) = \sum_{b=0}^{m-1} \left( i_b\, 2^{3b} + j_b\, 2^{3b+1} + k_b\, 2^{3b+2} \right).
\label{eq:morton}
\end{equation}
En lisant les bits de poids fort vers les poids faibles, la clé est la suite des triplets
$(k_b\, j_b\, i_b)$ pour $b = m-1, m-2, \dots, 0$ :
\[
K = \underbrace{k_{m-1}\, j_{m-1}\, i_{m-1}}_{\text{niveau } 1}\;
    \underbrace{k_{m-2}\, j_{m-2}\, i_{m-2}}_{\text{niveau } 2}\;\cdots\;
    \underbrace{k_0\, j_0\, i_0}_{\text{niveau } m}.
\]
Parcourir les cellules par clé croissante définit la \textbf{courbe de Morton}, ou
\emph{courbe en Z} (\emph{Z-order curve}), introduite par G.~M.~Morton en 1966 pour l'indexation
de données géographiques. C'est une courbe qui remplit l'espace : elle visite toutes les cellules,
en restant le plus souvent localisée.

La figure~\ref{fig:zorder} montre la version bidimensionnelle (deux bits par niveau,
$K = \sum_b i_b 2^{2b} + j_b 2^{2b+1}$) sur des grilles $4 \times 4$ et $8 \times 8$ : la courbe
parcourt le motif en \og Z\fg{} dans chaque bloc $2\times 2$, puis enchaîne les blocs eux-mêmes
selon le même motif, et ainsi de suite.

\begin{figure}[htb]
\centering
\begin{tikzpicture}[scale=0.9]
  % grille 4x4 avec numéros
  \draw[step=1,gray!50] (0,0) grid (4,4);
  \draw[thick] (0,0) rectangle (4,4);
  \draw[thick] (2,0) -- (2,4); \draw[thick] (0,2) -- (4,2);
  \draw[red!70!black,thick,->,>=stealth]
    (0.5,0.5) -- (1.5,0.5) -- (0.5,1.5) -- (1.5,1.5) -- (2.5,0.5) -- (3.5,0.5) -- (2.5,1.5) --
    (3.5,1.5) -- (0.5,2.5) -- (1.5,2.5) -- (0.5,3.5) -- (1.5,3.5) -- (2.5,2.5) -- (3.5,2.5) --
    (2.5,3.5) -- (3.5,3.5);
  \foreach \k/\x/\y in {0/0/0,1/1/0,2/0/1,3/1/1,4/2/0,5/3/0,6/2/1,7/3/1,8/0/2,9/1/2,10/0/3,
                        11/1/3,12/2/2,13/3/2,14/2/3,15/3/3}
    \node[font=\scriptsize,anchor=north west,inner sep=1pt,fill=white,fill opacity=0.8,text opacity=1] at (\x,\y+1) {\k};
  \node[below] at (2,-0.1) {$i$ (bit de poids faible)};
  \node[rotate=90,above] at (-0.1,2) {$j$};
  \begin{scope}[xshift=6cm,scale=0.5]
    \draw[step=1,gray!40] (0,0) grid (8,8);
    \draw[thick] (0,0) rectangle (8,8);
    \draw[thick] (4,0) -- (4,8); \draw[thick] (0,4) -- (8,4);
    \draw[red!70!black,thick]
      (0.5,0.5) -- (1.5,0.5) -- (0.5,1.5) -- (1.5,1.5) -- (2.5,0.5) -- (3.5,0.5) -- (2.5,1.5) --
      (3.5,1.5) -- (0.5,2.5) -- (1.5,2.5) -- (0.5,3.5) -- (1.5,3.5) -- (2.5,2.5) -- (3.5,2.5) --
      (2.5,3.5) -- (3.5,3.5) -- (4.5,0.5) -- (5.5,0.5) -- (4.5,1.5) -- (5.5,1.5) -- (6.5,0.5) --
      (7.5,0.5) -- (6.5,1.5) -- (7.5,1.5) -- (4.5,2.5) -- (5.5,2.5) -- (4.5,3.5) -- (5.5,3.5) --
      (6.5,2.5) -- (7.5,2.5) -- (6.5,3.5) -- (7.5,3.5) -- (0.5,4.5) -- (1.5,4.5) -- (0.5,5.5) --
      (1.5,5.5) -- (2.5,4.5) -- (3.5,4.5) -- (2.5,5.5) -- (3.5,5.5) -- (0.5,6.5) -- (1.5,6.5) --
      (0.5,7.5) -- (1.5,7.5) -- (2.5,6.5) -- (3.5,6.5) -- (2.5,7.5) -- (3.5,7.5) -- (4.5,4.5) --
      (5.5,4.5) -- (4.5,5.5) -- (5.5,5.5) -- (6.5,4.5) -- (7.5,4.5) -- (6.5,5.5) -- (7.5,5.5) --
      (4.5,6.5) -- (5.5,6.5) -- (4.5,7.5) -- (5.5,7.5) -- (6.5,6.5) -- (7.5,6.5) -- (6.5,7.5) --
      (7.5,7.5);
  \end{scope}
\end{tikzpicture}
\caption{Courbe de Morton en dimension 2. À gauche, grille $4\times4$ ($m = 2$) avec la clé de
chaque cellule ; à droite, grille $8\times8$ ($m = 3$). Les traits épais délimitent les cellules
du niveau 1 : chacune contient un intervalle contigu de clés.}
\label{fig:zorder}
\end{figure}

\subsection{La propriété fondamentale : préfixe = cellule parente}

Les $3$ bits de poids fort de la clé donnent la position de la cellule au niveau 1 (dans quel
octant de la boîte elle se trouve), les 6 bits de poids fort sa position au niveau 2, et plus
généralement :

\medskip
\noindent\fbox{\parbox{0.96\linewidth}{
\textbf{Propriété.} Soit une cellule de profondeur $m$ de clé $K$. La cellule de profondeur
$d \le m$ qui la contient a pour clé (dans la grille $2^d \times 2^d \times 2^d$)
\[
K_d = K \gg 3(m-d) = \left\lfloor \frac{K}{8^{\,m-d}} \right\rfloor,
\]
où $\gg$ désigne le décalage binaire à droite.}}
\medskip

\noindent\emph{Démonstration.} L'indice de la cellule ancêtre de profondeur $d$ dans la direction
$x$ est $\lfloor i / 2^{m-d} \rfloor$, dont les bits sont $i_{m-1}, \dots, i_{m-d}$ ; de même pour
$j$ et $k$. Entrelacer ces bits revient exactement à garder les $3d$ bits de poids fort de $K$,
c'est-à-dire à supprimer les $3(m-d)$ bits de poids faible. $\square$

\medskip
Deux conséquences rendent la clé de Morton parfaitement adaptée à notre problème :
\begin{enumerate}
  \item \textbf{Une seule clé suffit pour toutes les profondeurs} : on calcule la clé de chaque
        particule une fois, à la profondeur maximale, et on obtient la cellule à n'importe quelle
        profondeur par un simple décalage.
  \item \textbf{L'ordre de Morton est hiérarchique} : les cellules d'une même cellule parente
        ont des clés contiguës (figure~\ref{fig:zorder}). Une fois les clés triées, les
        particules d'une même cellule, à \emph{toute} profondeur, forment un bloc contigu du
        tableau trié. Compter les cellules non vides à la profondeur $d$ revient alors à compter
        les changements de valeur de $K \gg 3(m-d)$ le long du tableau trié.
\end{enumerate}

ScalFMM utilise lui-même les indices de Morton pour ranger les cellules
(\code{FTreeCoordinate::getMortonIndex}) ; c'est aussi le principe des octrees linéaires utilisés
par la plupart des codes FMM et \emph{treecode}.

\subsection{Exemple de comptage en dimension 2}

Prenons 8 particules dans une grille $4 \times 4$ ($m = 2$), situées dans les cellules
\[
(0,0),\ (0,0),\ (1,0),\ (1,1),\ (2,3),\ (3,3),\ (3,3),\ (2,0).
\]
Leurs clés (figure~\ref{fig:zorder}) sont $0, 0, 1, 3, 14, 15, 15, 4$, soit après tri :

\begin{center}
\begin{tabular}{lcccccccc}
\toprule
clé triée $K$ (profondeur 2)       & 0 & 0 & 1 & 3 & 4 & 14 & 15 & 15 \\
en binaire                         & 0000 & 0000 & 0001 & 0011 & 0100 & 1110 & 1111 & 1111 \\
$K \gg 2$ (profondeur 1)           & 0 & 0 & 0 & 0 & 1 & 3 & 3 & 3 \\
\bottomrule
\end{tabular}
\end{center}

\begin{itemize}
  \item À la profondeur 2, les valeurs distinctes sont $0, 1, 3, 4, 14, 15$ : 6 feuilles non
        vides, $8/6 \approx 1{,}3$ particule par feuille.
  \item À la profondeur 1, les valeurs distinctes sont $0, 1, 3$ : 3 cellules non vides sur 4
        (le quadrant $x < 2$, $y \ge 2$, de clé 2, est vide), $8/3 \approx 2{,}7$ particules par
        cellule.
\end{itemize}
Dans les deux cas, un seul parcours du tableau trié suffit : on incrémente le compteur à chaque
fois que la valeur décalée diffère de celle de l'élément précédent.

%======================================================================
\section{L'algorithme de choix de la hauteur}
%======================================================================

\subsection{Énoncé}

La règle, reprise de st-feeLLGood (commit \code{6063be5}, \code{src/fmm\_demag.cc}), est :

\medskip
\noindent\fbox{\parbox{0.96\linewidth}{
La hauteur de l'octree est le plus grand entier $h \in [h_{\min}, h_{\max}] = [3, 10]$ tel que les
feuilles non vides de l'octree de hauteur $h$ contiennent en moyenne au moins
$s_{\min} = 150$ particules :
\[
h = \max\left\{ h \in [4, 10] \;:\; N \ge s_{\min}\, L_h \right\},
\quad\text{ou } h = 3 \text{ si cet ensemble est vide.}
\]}}
\medskip

Le nombre $L_h$ de feuilles non vides croît avec $h$ (chaque feuille non vide de hauteur $h$
contient au moins une feuille non vide de hauteur $h+1$), donc $s_h = N/L_h$ décroît avec $h$.
L'ensemble des hauteurs admissibles est un intervalle $[4, h^\ast]$ et il suffit de parcourir les
hauteurs par valeurs décroissantes en s'arrêtant à la première qui satisfait le critère.

\subsection{Choix de la profondeur de référence}

Pour une hauteur $h$, les feuilles sont à la profondeur $h-1$. La plus grande profondeur utile
est donc $m = h_{\max} - 1 = 9$, soit une grille de $2^9 = 512$ cellules par direction. Chaque
indice tient sur 9 bits et la clé sur $3 \times 9 = 27$ bits, qui tiennent largement dans un
entier de 64 bits (\code{uint64\_t}). D'après la propriété du préfixe, la clé de la feuille
contenant la particule dans l'octree de hauteur $h$ est
\begin{equation}
K_{h-1} = K \gg 3\,(m - (h-1)) = K \gg 3\,(h_{\max} - h).
\label{eq:shift}
\end{equation}

\subsection{Quantification des positions}

Pour une particule de position normalisée $\tilde{\mathbf{x}}$, l'indice de la cellule de
profondeur $m$ dans la direction $d$ est
\begin{equation}
c_d = \min\left(\max\left(\left\lfloor \frac{\tilde{x}_d + w/2}{w}\, 2^m \right\rfloor,\, 0\right),\,
2^m - 1\right),
\label{eq:quant}
\end{equation}
où $\tilde{x}_d + w/2$ est la position relative au coin inférieur de la boîte. Le
\og clampage\fg{} entre $0$ et $2^m-1$ ne sert qu'aux particules exactement sur le bord supérieur
de la boîte, ou en dehors par erreur d'arrondi ; avec la marge de la boîte ($w = 2{,}01$ pour des
particules dans $[-1,1]^3$), il ne joue pas en pratique.

\paragraph{Cohérence avec ScalFMM.} ScalFMM calcule l'indice d'une particule dans la grille des
feuilles d'un octree de hauteur $h$ par
$\lfloor (\tilde{x}_d + w/2) / w_h \rfloor$ avec $w_h = w/2^{h-1}$
(\code{FCoordinateComputer::GetTreeCoordinate}). Or, pour un réel $t \ge 0$ et des entiers
$p \ge q$,
\[
\left\lfloor \frac{\lfloor 2^p\, t \rfloor}{2^{p-q}} \right\rfloor = \lfloor 2^q\, t \rfloor,
\]
donc décaler l'indice de profondeur $m$ de $m - (h-1)$ bits donne exactement l'indice de
ScalFMM (aux arrondis flottants près pour une particule à une frontière de cellule, ce qui ne
change pas le décompte de façon significative). Les feuilles que nous comptons sont donc
exactement celles que ScalFMM va créer.

L'entrelacement de ScalFMM place $z$ sur le bit de poids faible de chaque triplet, alors que le
nôtre (équation~\eqref{eq:morton}) y place $x$. Cela ne change rien : à chaque profondeur, les deux
conventions sont des bijections entre cellules et clés, et le nombre de valeurs distinctes est le
même.

\subsection{Exemple numérique en dimension 3}

Prenons la particule $\tilde{\mathbf{x}} = (0{,}5;\ -0{,}5;\ -0{,}5)$, avec $w = 2{,}01$ et
$2^m = 512$. L'équation~\eqref{eq:quant} donne
\begin{align*}
c_x &= \lfloor 1{,}505/2{,}01 \times 512 \rfloor = \lfloor 383{,}36 \rfloor = 383 = (101111111)_2,\\
c_y = c_z &= \lfloor 0{,}505/2{,}01 \times 512 \rfloor = \lfloor 128{,}64 \rfloor = 128 = (010000000)_2.
\end{align*}
En entrelaçant (triplets $z\,y\,x$ du poids fort au poids faible) :
\[
K = \underbrace{001}_{b=8}\ \underbrace{110}_{b=7}\ \underbrace{001}_{b=6}\
    \underbrace{001}_{b=5}\ \underbrace{001}_{b=4}\ \underbrace{001}_{b=3}\
    \underbrace{001}_{b=2}\ \underbrace{001}_{b=1}\ \underbrace{001}_{b=0}
  = 29\,659\,721.
\]
Les clés de feuilles s'obtiennent par décalage (équation~\eqref{eq:shift}) :
\begin{center}
\begin{tabular}{cccc>{\raggedright\arraybackslash}p{5cm}}
\toprule
$h$ & profondeur & décalage & clé de la feuille & interprétation \\
\midrule
2 & 1 & 24 & $(001)_2 = 1$ & octant $x > 0$, $y < 0$, $z < 0$ \\
3 & 2 & 21 & $(001\,110)_2 = 14$ & puis sous-octant $x$ bas, $y$ et $z$ hauts \\
4 & 3 & 18 & $(001\,110\,001)_2 = 113$ & \dots \\
\bottomrule
\end{tabular}
\end{center}
Ces valeurs sont celles calculées par la fonction \code{mortonKey} du code (vérifiées par un petit
programme de test).

\subsection{Algorithme complet}

\begin{algorithm}[htb]
\caption{Hauteur de l'octree adaptée aux particules (\code{octreeHeight})}
\label{alg:height}
\KwIn{positions normalisées $\tilde{\mathbf{x}}_1, \dots, \tilde{\mathbf{x}}_N$, largeur de boîte $w$}
\KwOut{hauteur $h$ de l'octree}
$m \leftarrow h_{\max} - 1$\;
\For{$p \leftarrow 1$ \KwTo $N$}{
  $K_p \leftarrow$ clé de Morton de $\tilde{\mathbf{x}}_p$ à la profondeur $m$
  (équations~\eqref{eq:quant} et~\eqref{eq:morton})\;
}
trier $K_1, \dots, K_N$ par ordre croissant\;
\For{$h \leftarrow h_{\max}$ \KwTo $h_{\min} + 1$ (par pas de $-1$)}{
  $\delta \leftarrow 3\,(h_{\max} - h)$\;
  $L \leftarrow 0$\;
  \For{$p \leftarrow 1$ \KwTo $N$}{
    \If{$p = 1$ \textbf{ou} $(K_p \gg \delta) \ne (K_{p-1} \gg \delta)$}{
      $L \leftarrow L + 1$ \tcp*{nouvelle feuille non vide}
    }
  }
  \If{$N \ge s_{\min}\, L$}{
    \KwRet{$h$}\;
  }
}
\KwRet{$h_{\min}$}\;
\end{algorithm}

\paragraph{Complexité.} Le calcul des clés est en $O(N\, m)$, le tri en $O(N \log N)$ et les
au plus $h_{\max} - h_{\min} = 7$ parcours en $O(N)$ chacun. On n'a jamais besoin de construire
d'octree, ni de stocker autre chose qu'un tableau de $N$ entiers de 64 bits (6 Mo pour
$7{,}5\times 10^5$ particules). Mesuré sur le portable (un seul thread, particules aléatoires) :
\begin{center}
\begin{tabular}{rcr}
\toprule
$N$ & hauteur obtenue & temps \\
\midrule
$110\,795$ & 4 & 17 ms \\
$656\,632$ & 5 & 92 ms \\
$7\,500\,000$ & 6 & 1,3 s \\
\bottomrule
\end{tabular}
\end{center}
Ce calcul n'est fait qu'une fois, à la construction de l'octree : son coût est celui de moins
d'un pas de temps et il est négligeable sur une simulation. Il pourrait être parallélisé
(tri parallèle) si de très gros maillages le rendaient gênant.

\subsection{\texorpdfstring{Les bornes $h_{\min} = 3$ et $h_{\max} = 10$}{Les bornes hmin = 3 et hmax = 10}}

\begin{itemize}
  \item $h_{\min} = 3$ : avec moins de 3 niveaux, il n'y a plus de liste d'interaction (au niveau
        1, les 8 cellules sont toutes voisines) et la FMM se réduit au calcul direct. Avec 3
        niveaux, les M2L n'ont lieu qu'au niveau 2 (64 cellules au plus) : c'est presque un
        calcul direct, ce qui convient aux petits maillages.
  \item $h_{\max} = 10$ : $8^9 \approx 1{,}3\times 10^8$ feuilles possibles ; avec 150 particules
        par feuille, il faudrait des milliards de particules pour l'atteindre (ou des
        particules extrêmement concentrées). La borne garantit aussi que la clé tient sur
        27 bits.
\end{itemize}

\subsection{\texorpdfstring{Calibration du seuil $s_{\min} = 150$}{Calibration du seuil smin = 150}}
\label{sec:calib}

Le seuil a été calibré dans st-feeLLGood, avec le même noyau (rotation, $P = 9$), sur un cylindre
court (\code{short\_cyl}) maillé de $2\times 10^4$ à $7{,}5\times 10^6$ particules, avec
8 threads, en comparant pour chaque maillage les temps obtenus avec plusieurs hauteurs. Conclusions :
\begin{itemize}
  \item l'octree le plus rapide a toujours environ 300 particules par feuille non vide :
        c'est la valeur $s^\ast$ de l'équation~\eqref{eq:sopt} pour ce noyau ;
  \item un niveau de plus (environ 35 particules par feuille) est 4 à 5 fois plus lent ;
  \item un niveau de moins (environ 2000 particules par feuille) n'est que 10 à 25~\% plus lent.
\end{itemize}

Comme passer d'une hauteur à la suivante divise $s_h$ par environ 8 (un peu moins si certaines
cellules filles sont vides), la règle \og plus grande hauteur avec $s_h \ge s_{\min}$\fg{} choisit
une hauteur telle que $s_{\min} \le s_h < 8\, s_{\min}$ environ, soit $s_h \in [150, 1200[$.
Cet intervalle contient l'optimum $s^\ast \approx 300$ et exclut le niveau de trop (35 particules
par feuille), qui est le plus pénalisant. Le seuil est volontairement placé du côté des
\emph{grandes} feuilles, puisqu'un niveau de moins coûte beaucoup moins cher qu'un niveau de trop.

\begin{center}
\begin{tabular}{rcrcr}
\toprule
particules & hauteur adaptée & temps FMM & ancienne hauteur (st) & ancien temps \\
\midrule
$19\,614$ & 3 & 7 ms & 8 & 260 ms \\
$1{,}3\times 10^5$ & 4 & 96 ms & 7 & 1,4 s \\
$10^6$ & 5 & 665 ms & 7 & 7,9 s \\
$7{,}5\times 10^6$ & 6 & 6,1 s & 7 & 26 s \\
\bottomrule
\end{tabular}

\smallskip
{\small Mesures de st-feeLLGood, 8 threads, temps d'un appel FMM.}
\end{center}

\subsection{\texorpdfstring{Le paramètre \code{SizeSubLevels}}{Le paramètre SizeSubLevels}}

ScalFMM découpe l'octree en sous-octrees de \code{SizeSubLevels} niveaux ; chaque sous-octree
alloue des tableaux de pointeurs de taille $8^\ell$ pour ses niveaux $\ell$ (les cellules
elles-mêmes ne sont créées que si elles sont non vides). ScalFMM exige
$\code{SizeSubLevels} \le h - 1$ (assertion dans le constructeur de \code{FOctree}). On garde la
valeur 3 utilisée jusqu'ici, réduite à $h - 1$ si nécessaire : $\min(3, h-1)$, comme dans
st-feeLLGood. Pour $h = 3$, cela donne 2.

%======================================================================
\section{Mise en œuvre dans feeLLGood}
%======================================================================

La mise en œuvre se trouve dans la branche \code{optim\_tree\_height}, en deux commits :
\begin{itemize}
  \item \code{779944b} \emph{Adapt the FMM octree height to the particles} ;
  \item \code{f8e0a7f} \emph{Build the FMM algorithm once} ;
\end{itemize}
suivis d'un commit de documentation (\code{f42210d}).

\subsection{\texorpdfstring{Le fichier \code{src/octree\_height.h}}{Le fichier src/octree\_height.h}}

Les fonctions de calcul de la hauteur sont dans un en-tête séparé, sans dépendance à ScalFMM
(seulement la bibliothèque standard et Eigen). Cela permet de les tester unitairement sans
compiler ScalFMM, et de les réutiliser telles quelles si feeLLGood passe à ScalFMM~3. Elles sont
dans l'espace de noms \code{scal\_fmm}, comme le reste de l'interface FMM.

\begin{lstlisting}[caption={Constantes et clé de Morton (\code{src/octree\_height.h}).}]
const int MinLevels = 3;                /**< minimal height of the octree */
const int MaxLevels = 10;               /**< maximal height of the octree */
const double MinParticlesPerLeaf = 150; /**< minimal mean number of particles per non-empty leaf */

inline uint64_t mortonKey(const Eigen::Ref<const Eigen::Vector3d> p,
                          const double boxWidth)
    {
    const int n = 1 << (MaxLevels - 1);
    uint64_t key = 0;
    for (int d = 0; d < 3; d++)
        {
        int c = static_cast<int>((p(d) + boxWidth / 2) / boxWidth * n);
        c = std::clamp(c, 0, n - 1);
        for (int b = 0; b < MaxLevels - 1; b++)
            { key |= uint64_t((c >> b) & 1) << (3 * b + d); }
        }
    return key;
    }
\end{lstlisting}

\begin{itemize}
  \item \code{n} $= 2^{h_{\max}-1} = 512$ est le nombre de cellules par direction à la profondeur
        de référence.
  \item Ligne 12 : quantification~\eqref{eq:quant}. La conversion \code{static\_cast<int>}
        tronque vers zéro, ce qui coïncide avec la partie entière pour un argument positif ; un
        argument négatif (particule hors de la boîte) donne $0$ ou une valeur négative ramenée à
        $0$ par \code{std::clamp}.
  \item Lignes 14--15 : entrelacement~\eqref{eq:morton}. Le bit $b$ de l'indice de la direction
        $d$ ($d = 0, 1, 2$ pour $x, y, z$) est placé en position $3b + d$ de la clé. La
        conversion en \code{uint64\_t} \emph{avant} le décalage est nécessaire : décaler un
        \code{int} de plus de 31 bits serait un comportement indéfini (ici les positions vont
        jusqu'à 26, mais la conversion rend le code sûr pour d'autres valeurs de
        \code{MaxLevels}, jusqu'à 22).
\end{itemize}

\begin{lstlisting}[caption={Calcul de la hauteur (\code{src/octree\_height.h}).}]
inline int octreeHeight(const std::vector<Eigen::Vector3d> &pos,
                        const double boxWidth)
    {
    std::vector<uint64_t> keys(pos.size());
    std::transform(pos.begin(), pos.end(), keys.begin(),
                   [boxWidth](const Eigen::Vector3d &p) { return mortonKey(p, boxWidth); });
    std::sort(keys.begin(), keys.end());
    for (int h = MaxLevels; h > MinLevels; h--)
        {
        const int shift = 3 * (MaxLevels - h);
        size_t nbLeaves = 0;
        for (size_t i = 0; i < keys.size(); i++)
            {
            if (i == 0 || (keys[i] >> shift) != (keys[i - 1] >> shift))
                { nbLeaves++; }
            }
        if (double(keys.size()) >= MinParticlesPerLeaf * double(nbLeaves))
            { return h; }
        }
    return MinLevels;
    }
\end{lstlisting}

C'est la traduction directe de l'algorithme~\ref{alg:height}. Le test est écrit
$N \ge s_{\min} L_h$ plutôt que $N / L_h \ge s_{\min}$ pour éviter une division. Remarque : pour
un ensemble vide de particules, $L_h = 0$ et le test est vrai dès $h = 10$ ; ce cas ne se
présente pas (un maillage sans élément magnétique n'a pas de champ démagnétisant à calculer).

\subsection{\texorpdfstring{Modifications de \code{src/fmm\_demag.h}}{Modifications de src/fmm\_demag.h}}

La classe \code{scal\_fmm::fmm} contient l'octree (\code{OctreeClass tree}) et le noyau
(\code{KernelClass kernels}) comme membres, construits dans la liste d'initialisation du
constructeur. Leur hauteur doit donc être connue \emph{avant} leur construction. En C++, les
membres sont initialisés dans l'ordre de leur \emph{déclaration} dans la classe (et non dans
l'ordre de la liste d'initialisation). On a donc :
\begin{enumerate}
  \item déplacé la déclaration du coefficient de normalisation \code{norm} avant celles de
        \code{tree} et \code{kernels}, et l'avons rendu constant, initialisé dans la liste
        d'initialisation (il était auparavant affecté dans le corps du constructeur) ;
  \item ajouté un membre constant \code{nbLevels}, déclaré entre \code{norm} et \code{tree} ;
  \item supprimé la constante globale \code{NbLevels = 6}.
\end{enumerate}

\begin{lstlisting}[caption={Membres et liste d'initialisation (\code{src/fmm\_demag.h}).}]
    // liste d'initialisation du constructeur
        : prmTetra(prmTet), prmTriangle(prmTri), norm(2. / msh.l.maxCoeff()),
          nbLevels(octreeHeight(particlePositions(msh), boxWidth)),
          tree(nbLevels, std::min(SizeSubLevels, nbLevels - 1), boxWidth, boxCenter),
          kernels(nbLevels, boxWidth, boxCenter)

    // membres prives, dans l'ordre de declaration
    const double norm;      // normalization coefficient
    const int nbLevels;     // height of the octree, adapted to the particles
    OctreeClass tree;       // tree initialized by constructor
    KernelClass kernels;    // kernel initialized by constructor
\end{lstlisting}

Les positions des particules sont fournies par une méthode privée \code{particlePositions}, qui
reproduit exactement l'ordre et la normalisation de l'insertion dans l'octree : nœuds magnétiques
(cibles), puis points de Gauss des tétraèdres magnétiques, puis points de Gauss des triangles
magnétiques (sources). Elle n'utilise que \code{norm} et le maillage, déjà disponibles au moment
de l'initialisation de \code{nbLevels}.

\begin{lstlisting}[caption={Positions normalisées des particules (\code{src/fmm\_demag.h}).}]
std::vector<Eigen::Vector3d> particlePositions(const Mesh::mesh &msh) const
    {
    std::vector<Eigen::Vector3d> pos;
    pos.reserve(msh.magNode.size() + msh.magTet.size() * Tetra::NPI
                + msh.magTri.size() * Triangle::NPI);
    for (size_t i = 0; i < msh.magNode.size(); i++)
        {
        if (msh.magNode[i])
            { pos.push_back(norm * (msh.getNode_p(i) - msh.c)); }
        }
    for (const int idx : msh.magTet)
        {
        Eigen::Matrix<double, Nodes::DIM, Tetra::NPI> gauss = msh.tet[idx].getPtGauss();
        for (int j = 0; j < Tetra::NPI; j++)
            { pos.push_back(norm * (gauss.col(j) - msh.c)); }
        }
    for (const int idx : msh.magTri)
        {
        Eigen::Matrix<double, Nodes::DIM, Triangle::NPI> gauss = msh.tri[idx].getPtGauss();
        for (int j = 0; j < Triangle::NPI; j++)
            { pos.push_back(norm * (gauss.col(j) - msh.c)); }
        }
    return pos;
    }
\end{lstlisting}

Le code d'insertion dans l'octree n'a pas été modifié : il recalcule les mêmes positions. Ce
double calcul, fait une seule fois au démarrage, coûte quelques millisecondes ; il évite de
toucher à l'insertion, dont les indices de particules (\code{idxPart}, qui parcourt aussi les
nœuds non magnétiques) servent ensuite à retrouver les densités de charge.

Enfin, une méthode publique \code{getNbLevels()} donne la hauteur choisie, affichée en mode
verbeux par \code{src/main.cpp} :
\begin{lstlisting}[language={},numbers=none,frame=none,basicstyle=\ttfamily\footnotesize]
Magnetostatics: particles inserted, using 8 threads, octree of 4 levels, in 28.21 ms
\end{lstlisting}

\subsection{L'algorithme FMM construit une seule fois}

Le second commit reprend l'autre optimisation du commit de st-feeLLGood. L'objet algorithme de
ScalFMM (\code{FmmClass}, c'est-à-dire \code{FFmmAlgorithmThreadTsm}) était construit à chaque
appel de \code{fmm::demag}, donc deux fois par pas de temps. Or son constructeur
\begin{itemize}
  \item lit le nombre de threads OpenMP (\code{omp\_get\_max\_threads()}) et le fige ;
  \item fait une copie du noyau pour chaque thread, avec ses tables de rotation précalculées.
\end{itemize}
Il est maintenant construit une seule fois, à la fin du constructeur de \code{scal\_fmm::fmm},
et conservé dans un membre \code{std::unique\_ptr<FmmClass> algo}. Il est créé \emph{après}
l'appel à \code{omp\_set\_num\_threads(ScalfmmNbThreads)}, de sorte que le nombre de threads
utilisé par la FMM est inchangé. On ne peut pas en faire un membre ordinaire initialisé dans la
liste d'initialisation, puisque celle-ci est exécutée avant le corps du constructeur, donc avant
\code{omp\_set\_num\_threads}.

\begin{lstlisting}[caption={Construction unique de l'algorithme (\code{src/fmm\_demag.h}).}]
        // fin du constructeur
        // after omp_set_num_threads(): it copies the kernel for each thread
        algo = std::make_unique<FmmClass>(&tree, &kernels);

    // dans fmm::demag, a chaque appel
        algo->execute();
\end{lstlisting}

\subsection{Tests unitaires}

Le fichier \code{unit-tests/ut\_octree\_height.cpp} (cible \code{test\_ut\_octree\_height}) teste
les deux fonctions sur des configurations dont le résultat se calcule à la main. La fonction
\code{grid(n, a, w)} y place $n^3$ particules aux centres des cellules d'une grille régulière qui
remplit le cube $[a, a+w]^3$.

\begin{description}
  \item[\code{morton\_key\_parent}] Avec un décalage de $3(h_{\max} - 2)$, la clé de profondeur 1
        doit valoir $i + 2j + 4k$ où $(i, j, k)$ repère l'octant : on vérifie les points
        $(\pm 0{,}5, \pm 0{,}5, \pm 0{,}5)$ (clés 0, 1, 2, 4 et 7). Un point hors de la boîte,
        $(-2, -2, -2)$, doit avoir la clé 0 (clampage).
  \item[\code{min\_levels}] $4^3 = 64$ particules : $64 < 150$ même pour une seule feuille, d'où
        $h = 3$.
  \item[\code{uniform\_cube}] $64^3 = 262\,144$ particules remplissant la boîte. À la hauteur
        $h$, il y a $8^{h-1}$ feuilles non vides contenant chacune $262\,144/8^{h-1}$
        particules : 512 pour $h = 4$ ($\ge 150$), 64 pour $h = 5$ ($< 150$). D'où $h = 4$.
  \item[\code{empty\_leaves\_ignored}] Les mêmes $64^3$ particules dans le cube $[0, w/16]^3$,
        qui est exactement une cellule de profondeur 4. Les feuilles non vides de profondeur
        $h-1 \ge 4$ sont au nombre de $8^{h-5}$ : 512 particules par feuille pour $h = 8$,
        64 pour $h = 9$. D'où $h = 8$, soit 4 niveaux de plus que pour le cube plein : le
        test vérifie que seules les feuilles non vides sont comptées.
  \item[\code{max\_levels}] 1000 particules au même point : une seule feuille non vide à toutes
        les hauteurs, d'où $h = h_{\max} = 10$.
\end{description}

\subsection{Validation}

Changer la hauteur de l'octree change l'approximation FMM, mais au niveau de sa propre erreur
(de l'ordre de $10^{-4}$ en relatif sur le potentiel pour le noyau rotation à $P = 9$, nettement
moins sur les grandeurs intégrées) :
\begin{itemize}
  \item \code{ci-tests/full\_test.py} (ellipsoïde, 20 ps) passe, avec un écart de $10^{-6}$ sur
        l'aimantation moyenne et $6{,}5\times 10^{-7}$ en relatif sur l'énergie totale, pour un
        seuil de $10^{-5}$ ;
  \item sur \code{example0.1}, l'écart relatif de l'énergie totale avec master après 10 ps est
        de $2{,}5\times 10^{-7}$ ;
  \item la construction unique de l'algorithme ne change pas les résultats (écarts relatifs de
        $10^{-13}$, dus à l'ordre des sommes parallèles).
\end{itemize}

%======================================================================
\section{Résultats}
\label{sec:results}
%======================================================================

Les mesures portent sur trois simulations d'ellipsoïdes existantes, avec des durées simulées
courtes ($10^{-11}$ ou $2\times 10^{-12}$~s). Le \og temps magnétostatique par pas\fg{} est le
temps des deux appels FMM d'un pas de temps (potentiels de $\mathbf{u}$ et de $\mathbf{v}$), relevé
dans la sortie verbeuse (\code{magnetostatics done in \dots}) et moyenné sur tous les pas.

\subsection{Portable (8 threads)}

Intel Core i7-8650U, 8 threads, WSL2. Comparaison de master (hauteur 6) et de la hauteur
adaptée (premier commit seul), tableau~\ref{tab:laptop}.

\begin{table}[htbp]
\centering
\caption{Portable (i7-8650U, 8 threads) : master ($h = 6$) et hauteur adaptée (premier commit).}
\label{tab:laptop}
\begin{tabular}{lrrcrrrrr}
\toprule
maillage & nœuds & particules & $h$ & \multicolumn{2}{c}{magnéto. par pas} &
\multicolumn{2}{c}{temps par pas} & gain \\
\cmidrule(lr){5-6}\cmidrule(lr){7-8}
 & & & & master & adaptée & master & adaptée & \\
\midrule
\code{example0.1}        & 3990  & 110\,795 & 4 & 1130 ms & 169 ms  & 1,34 s & 0,35 s & 3,8 \\
\code{fmr\_5587nodes}    & 5587  & 158\,571 & 4 & 2057 ms & 330 ms  & 2,18 s & 0,42 s & 5,2 \\
\code{enhanced\_curling} & 22061 & 656\,632 & 5 & 4214 ms & 1148 ms & 4,57 s & 1,46 s & 3,1 \\
\bottomrule
\end{tabular}
\end{table}

L'effet de la construction unique de l'algorithme (second commit) a été mesuré sur la même
machine, avec deux exécutions de chaque version (tableau~\ref{tab:algo}).
\begin{table}[htbp]
\centering
\caption{Effet de la construction unique de l'algorithme FMM : temps magnétostatique par pas (portable, 8 threads).}
\label{tab:algo}
\begin{tabular}{lcc}
\toprule
maillage & algorithme construit à chaque appel & algorithme construit une fois \\
\midrule
\code{example0.1}     & 178 à 202 ms & 147 à 148 ms \\
\code{fmr\_5587nodes} & 408 à 418 ms & 221 à 231 ms \\
\bottomrule
\end{tabular}
\end{table}

\subsection{Machine de 96 cœurs (jusqu'à 192 threads)}

AMD EPYC 9654 (96 cœurs, 384 threads matériels). Master et la branche complète (deux commits),
compilés avec la même configuration (TBB, MKL), $2\times 10^{-12}$~s simulées. Temps
magnétostatique par pas en fonction du nombre de threads de la FMM
(\code{demagnetizing\_field\_solver: nb\_threads}) : tableau~\ref{tab:rigel}.

\begin{table}[htbp]
\centering
\caption{AMD EPYC 9654 : temps magnétostatique par pas selon le nombre de threads FMM.}
\label{tab:rigel}
\begin{tabular}{lcrrr}
\toprule
maillage (particules), hauteur adaptée & threads & master ($h = 6$) & branche & gain \\
\midrule
\code{example0.1} (110\,795), $h = 4$ & 8   & 152 ms & 37 ms & 4,1 \\
                                       & 32  & 97 ms  & 16 ms & 6,0 \\
                                       & 96  & 165 ms & 13 ms & 12,6 \\
                                       & 192 & 278 ms & 15 ms & 18,5 \\
\midrule
\code{fmr\_5587nodes} (158\,571), $h = 4$ & 8   & 229 ms & 50 ms & 4,6 \\
                                          & 32  & 126 ms & 20 ms & 6,4 \\
                                          & 96  & 178 ms & 18 ms & 10,1 \\
                                          & 192 & 322 ms & 19 ms & 16,9 \\
\midrule
\code{enhanced\_curling} (656\,632), $h = 5$ & 8   & 475 ms & 179 ms & 2,6 \\
                                             & 32  & 215 ms & 73 ms  & 2,9 \\
                                             & 96  & 239 ms & 53 ms  & 4,5 \\
                                             & 192 & 352 ms & 55 ms  & 6,4 \\
\bottomrule
\end{tabular}
\end{table}

Avec 6 niveaux, master \emph{ralentit} au-delà de 32 threads : les nombreuses petites cellules
des niveaux profonds rendent dominantes les synchronisations entre threads, à chaque niveau et
chaque passe. Avec la hauteur adaptée, le temps sature entre 32 et 96 threads pour ces tailles de
maillage, et ne se dégrade pas au-delà.

\subsection{Vérification du seuil : balayage des hauteurs}

Pour vérifier que le seuil de 150, calibré à 8 threads, reste pertinent sur une machine
multicœur, la hauteur a été forcée de 3 à 7 (par une variable d'environnement ajoutée
temporairement, qui ne fait pas partie du code) : tableau~\ref{tab:scan}.

\begin{table}[htbp]
\centering
\caption{AMD EPYC 9654 : temps magnétostatique par pas (ms) pour des hauteurs forcées de 3 à 7 ; en gras, la hauteur la plus rapide.}
\label{tab:scan}
\begin{tabular}{lcrrrrrc}
\toprule
maillage & threads & $h = 3$ & $h = 4$ & $h = 5$ & $h = 6$ & $h = 7$ & hauteur adaptée \\
\midrule
\code{example0.1}        & 8   & 125  & \textbf{37} & 39           & 137 & 296  & 4 \\
                         & 192 & 91   & \textbf{15} & 18           & 39  & 104  & 4 \\
\code{fmr\_5587nodes}    & 8   & 163  & \textbf{50} & 62           & 221 & 464  & 4 \\
                         & 192 & 108  & \textbf{18} & 24           & 63  & 158  & 4 \\
\code{enhanced\_curling} & 8   & 2097 & 443         & \textbf{182} & 460 & 1229 & 5 \\
                         & 192 & 940  & 85          & \textbf{54}  & 107 & 303  & 5 \\
\bottomrule
\end{tabular}
\end{table}

La hauteur adaptée est toujours la plus rapide, à 8 comme à 192 threads. Un niveau d'écart avec
l'optimum coûte de 1,1 à 3 fois plus, deux niveaux jusqu'à 10 fois plus. On retrouve
l'allure prédite par le modèle~\eqref{eq:cost} : une vallée étroite autour de l'optimum, avec une
remontée rapide des deux côtés. À 192 threads, les hauteurs 4 et 5 se rapprochent sur les petits
maillages (plus de cellules donnent plus de parallélisme), mais l'ordre ne s'inverse pas.

%======================================================================
\section{Conclusion et perspectives}
%======================================================================

La hauteur de l'octree est un paramètre déterminant pour le coût de la FMM, et l'optimum dépend
du nombre de particules par feuille non vide, pas de la taille du maillage. La règle reprise de
st-feeLLGood (la plus grande hauteur avec au moins 150 particules par feuille non vide en moyenne)
se calcule simplement et à faible coût grâce aux clés de Morton : une clé par particule, un tri,
puis un comptage des préfixes distincts pour chaque hauteur candidate. Elle choisit la hauteur
optimale sur tous les cas mesurés, sur un portable comme sur une machine de 96 cœurs, et divise le
temps magnétostatique par 2,6 à 18,5 sans perte de précision.

Pistes possibles :
\begin{itemize}
  \item vérifier le seuil sur de plus gros maillages ($10^7$ particules et plus) sur la machine
        multicœur, et pour d'autres ordres $P$ (l'optimum $s^\ast$ croît comme $P^{3/2}$ d'après
        l'équation~\eqref{eq:sopt}) ;
  \item ajouter un paramètre permettant d'imposer la hauteur, pour les tests et les études de
        performance ;
  \item paralléliser le tri des clés si le calcul de la hauteur devenait sensible sur de très
        gros maillages ;
  \item lors du passage à ScalFMM~3, réutiliser \code{octree\_height.h} tel quel : la
        fonction ne dépend que des positions et de la largeur de la boîte.
\end{itemize}

\end{document}
